\documentclass[11pt,a4paper]{article}
\usepackage[margin=2.1cm]{geometry}
\usepackage{amsmath,amssymb,booktabs,longtable,array,hyperref,xcolor}
\hypersetup{colorlinks=true,linkcolor=blue,urlcolor=blue}
\newcommand{\code}[1]{\texttt{#1}}
\title{A056 v0.3.0\\PKLSA Spectral--Sine-Gordon--Mittag-Leffler Falsifier}
\author{SST-Workbench}
\date{Framework v1.0.4 CANONICAL\_FROZEN}
\begin{document}\maketitle

% SST-REPORT-SECTION:IDENTIFICATION
\section{Identification}
Catalog A056, release v0.3.0. This release adds the real E010/PKLSA finite-core/Euler dynamic-provider route. This falsifier is a thin scientific instance of the immutable SST Falsifier Framework v1.0.4. The framework is referenced through \code{.sst\_framework\_root}; it is not vendored or patched by this release.

% SST-REPORT-SECTION:QUESTION
\section{Question}
For source-native E010/PKLSA-qualified vortex carriers evolved by a preregistered finite-core filament or incompressible Euler provider, do the extracted paired Kelvin residual phase and mode envelope require Sine--Gordon restoring dynamics and/or Mittag--Leffler fractional-memory relaxation?

% SST-REPORT-SECTION:HYPOTHESES
\section{Hypotheses}
The phase null family consists of Wave, Klein--Gordon and Duffing laws. The phase candidate is Sine--Gordon. The relaxation null family consists of exponential, stretched-exponential and bi-exponential laws; Mittag--Leffler is the fractional-memory candidate. A physical conclusion additionally requires canonical-framework eligible cross-source replication.

% SST-REPORT-SECTION:ASSUMPTIONS
\section{Assumptions and interpretation boundary}
For real E010 campaigns the phase observable is frozen as the paired Kelvin residual phase before scoring. Time and arclength-like coordinates are monotone and approximately uniform. Python/NumPy FP64 is authoritative reference arithmetic; C++/OpenMP FP64 certifies the A056 kernels. SYCL FP32 and DD32/FP32x2 remain screening-only. In conservative inviscid dynamics, projected ringdown is not automatically interpreted as thermodynamic dissipation.

% SST-REPORT-SECTION:SYMBOLS
\section{Symbols}
\begin{longtable}{lll}\toprule Symbol & Meaning & Units\\\midrule
$\varphi(s,t)$ & provider-declared phase/torsion observable & rad / dimensionless\\
$a$ & phase-wave coefficient & $L^2T^{-2}$\\
$b,c$ & restoring coefficients & $T^{-2}$ for dimensionless phase\\
$\alpha$ & Mittag--Leffler fractional order & 1\\
$\tau$ & relaxation timescale & $T$\\\bottomrule\end{longtable}

% SST-REPORT-SECTION:EQUATIONS
\section{Equations}
\subsection{Phase family}
\begin{align}
H_W &: \varphi_{tt}=a\varphi_{ss},\\
H_{KG} &: \varphi_{tt}=a\varphi_{ss}-b\varphi,\\
H_D &: \varphi_{tt}=a\varphi_{ss}-b\varphi-c\varphi^3,\\
H_{SG} &: \varphi_{tt}=a\varphi_{ss}-b\sin\varphi.
\end{align}
For dimensionless phase, $[a]=L^2T^{-2}$ and $[b]=[c]=T^{-2}$.

\subsection{Relaxation family}
\begin{align}
R_{\exp}(t)&=e^{-t/\tau},\\
R_{SE}(t)&=e^{-(t/\tau)^\beta},\\
R_{2e}(t)&=w e^{-t/\tau_1}+(1-w)e^{-t/\tau_2},\\
R_{ML}(t)&=E_{\alpha,1}[-(t/\tau)^\alpha],\\
E_{\alpha,\beta}(z)&=\sum_{k=0}^{\infty}\frac{z^k}{\Gamma(\alpha k+\beta)}.
\end{align}
The exact null limit $E_{1,1}(-t/\tau)=e^{-t/\tau}$ is explicitly protected by the registered interior support interval $0.48<\alpha<0.97$.

% SST-REPORT-SECTION:ALGORITHM
\section{Algorithm}
The v0.3.0 provider first resolves and re-hashes E010 source-native carriers, creates paired base/perturbed finite-core evolutions, advects material markers, and extracts the frozen phase/envelope. The scoring pipeline then validates paired NPZ/JSON provider cases, computes centered POD/SVD qualification, forms second-order temporal and spatial phase derivatives, fits all registered model families on discovery blocks, evaluates untouched confirmation blocks, performs deterministic twofold coarsening, and certifies the experiment-specific kernels against C++/OpenMP FP64. Replication is then assessed with the frozen framework evidence classes.

\IfFileExists{AUTO_SCIENCE_CONTRACT.tex}{% AUTO-GENERATED. Do not edit; generated from frozen science_contract.json.
\subsection*{Machine-readable science contract}
\begin{description}
\item[Research question] For source-native E010/PKLSA-qualified vortex carriers evolved by a preregistered finite-core filament or incompressible Euler provider, do the extracted paired Kelvin residual phase and mode envelope require Sine-Gordon restoring dynamics and/or Mittag-Leffler fractional-memory relaxation?
\item[Objective] Close the v0.2.3 provider gap by generating phi(s,t) and R(t) directly from qualified E010 carriers under frozen finite-core/Euler dynamics, then apply the canonical A056 blind competition without promoting simulation evidence to physical cross-source support.
\item[$H_0$] H0: Sine-Gordon and Mittag-Leffler descriptions do not achieve the preregistered discovery, held-out, numerical-certification and independent-replication thresholds beyond the registered alternatives.
\item[$H_1$] H1: one or both candidate mechanisms outperform the registered alternatives, survive held-out and numerical certification, and replicate across eligible independent physical evidence.
\end{description}
\subsubsection*{Registered equations}
\paragraph{E1: Wave null phase model}
\begin{equation}
\varphi_{tt}=a\varphi_{ss}
\end{equation}
\noindent\textit{Dimensional check:} Both sides have dimensions T\textasciicircum{}\{-2\} for dimensionless phase when [a]=L\textasciicircum{}2 T\textasciicircum{}\{-2\}.\\
\paragraph{E2: Klein-Gordon competitor}
\begin{equation}
\varphi_{tt}=a\varphi_{ss}-b\varphi
\end{equation}
\noindent\textit{Dimensional check:} [b]=T\textasciicircum{}\{-2\}; every term has dimensions T\textasciicircum{}\{-2\}.\\
\paragraph{E3: Duffing competitor}
\begin{equation}
\varphi_{tt}=a\varphi_{ss}-b\varphi-c\varphi^3
\end{equation}
\noindent\textit{Dimensional check:} For dimensionless phase, [b]=[c]=T\textasciicircum{}\{-2\}.\\
\paragraph{E4: Sine-Gordon candidate}
\begin{equation}
\varphi_{tt}=a\varphi_{ss}-b\sin\varphi
\end{equation}
\noindent\textit{Dimensional check:} sin(varphi) is dimensionless and [b]=T\textasciicircum{}\{-2\}.\\
\paragraph{E5: Bayesian information criterion used on discovery blocks}
\begin{equation}
\mathrm{BIC}=n\ln(\mathrm{RSS}/n)+k\ln n
\end{equation}
\noindent\textit{Dimensional check:} Model residuals are compared within one observable/scaling; BIC differences are dimensionless.\\
\paragraph{E6: Held-out normalized root-mean-square error}
\begin{equation}
\mathrm{NRMSE}=\sqrt{n^{-1}\sum_i(y_i-\hat y_i)^2}/\sigma_y
\end{equation}
\noindent\textit{Dimensional check:} Numerator and denominator have the same observable units, so NRMSE is dimensionless.\\
\paragraph{E7: Mittag-Leffler function}
\begin{equation}
E_{\alpha,\beta}(z)=\sum_{k=0}^{\infty} z^k/\Gamma(\alpha k+\beta)
\end{equation}
\noindent\textit{Dimensional check:} z must be dimensionless.\\
\paragraph{E8: Mittag-Leffler relaxation candidate}
\begin{equation}
R_{ML}(t)=E_{\alpha,1}[-(t/\tau)^\alpha]
\end{equation}
\noindent\textit{Dimensional check:} t/tau is dimensionless.\\
\paragraph{E9: Ordinary exponential null limit}
\begin{equation}
R_{exp}(t)=\exp(-t/\tau)=E_{1,1}(-t/\tau)
\end{equation}
\noindent\textit{Dimensional check:} t/tau is dimensionless.\\
\paragraph{E10: Stretched-exponential competitor}
\begin{equation}
R_{SE}(t)=\exp[-(t/\tau)^\beta]
\end{equation}
\noindent\textit{Dimensional check:} t/tau is dimensionless.\\
\paragraph{E11: Bi-exponential competitor}
\begin{equation}
R_{2e}(t)=w e^{-t/\tau_1}+(1-w)e^{-t/\tau_2}
\end{equation}
\noindent\textit{Dimensional check:} Each exponential argument is dimensionless and w is dimensionless.\\
\paragraph{E12: Phase-resolution convergence diagnostic}
\begin{equation}
\delta_c=\lVert\mathbf c_f-\mathbf c_c\rVert_2/\max(\lVert\mathbf c_f\rVert_2,10^{-30}),\quad \mathbf c=(a,b)
\end{equation}
\noindent\textit{Dimensional check:} Coefficient vectors are compared componentwise under the same coordinate units; ratio is dimensionless.\\
\paragraph{E13: Memory-resolution convergence diagnostics}
\begin{equation}
\delta_\alpha=|\alpha_f-\alpha_c|,\qquad \delta_\tau=|\tau_f-\tau_c|/\max(|\tau_f|,10^{-30})
\end{equation}
\noindent\textit{Dimensional check:} Both diagnostics are dimensionless.\\
\paragraph{E14: Incompressible unforced Euler evolution}
\begin{equation}
\partial_t\mathbf u+(\mathbf u\cdot\nabla)\mathbf u=-\nabla p,\qquad \nabla\cdot\mathbf u=0
\end{equation}
\noindent\textit{Dimensional check:} Each acceleration term has dimensions L T\textasciicircum{}\{-2\}; incompressibility has T\textasciicircum{}\{-1\}.\\
\paragraph{E15: Finite-core centerline vorticity seed}
\begin{equation}
\boldsymbol\omega(\mathbf x,0)\propto\sum_j \mathbf T_j\exp[-d_{per}(\mathbf x,\mathbf X_j)^2/(2\sigma^2)]
\end{equation}
\noindent\textit{Dimensional check:} The exponential is dimensionless; normalization fixes the velocity scale.\\
\paragraph{E16: Material-marker transport}
\begin{equation}
\dot{\mathbf X}(s,t)=\mathbf u(\mathbf X(s,t),t)
\end{equation}
\noindent\textit{Dimensional check:} Both sides have dimensions L T\textasciicircum{}\{-1\}.\\
\paragraph{E17: Frozen Kelvin perturbation}
\begin{equation}
\mathbf X_\epsilon(s,0)=\mathbf X_0(s)+\epsilon[\cos(m\theta)\mathbf e_1+\sin(m\theta)\mathbf e_2]
\end{equation}
\noindent\textit{Dimensional check:} epsilon has dimensions L; trigonometric factors and frame vectors are dimensionless.\\
\paragraph{E18: Complex paired transverse displacement}
\begin{equation}
q(s,t)=\delta\mathbf X\cdot\mathbf e_1+i\,\delta\mathbf X\cdot\mathbf e_2
\end{equation}
\noindent\textit{Dimensional check:} q has dimensions L.\\
\paragraph{E19: Residual Kelvin phase}
\begin{equation}
\varphi(s,t)=\arg[q(s,t)e^{-im\theta(s)}]
\end{equation}
\noindent\textit{Dimensional check:} The argument is an angle and therefore dimensionless.\\
\paragraph{E20: Projected mode envelope}
\begin{equation}
R(t)=|c_m(t)|/|c_m(0)|,\qquad c_m=N_s^{-1}\sum_j q(s_j,t)e^{-im\theta_j}
\end{equation}
\noindent\textit{Dimensional check:} R is dimensionless.\\
\subsubsection*{Registered code-to-mathematics steps}
\begin{enumerate}
\item \textbf{S1} [G0; equations E1]: Verify frozen protocol/blind commitments and enumerate hashed provider cases.
\item \textbf{S0P} [G0; equations E14, E15, E16, E17, E18, E19, E20]: Resolve and re-hash E010 v0.3.1 strict upstream-provider representatives, construct finite-core base/perturbed carriers, evolve them with the frozen provider, and commit the private carrier reveal before A056 scoring.
\item \textbf{S2} [G1; equations E1]: Validate provider metadata, array shapes, monotonicity, finite values and sampling uniformity.
\item \textbf{S3} [G2; equations E4, E9]: Evaluate Python FP64 finite differences, POD spectral stability and the Mittag-Leffler alpha=1 exponential identity.
\item \textbf{S4} [G3; equations E1, E2, E3, E4, E5, E8, E9, E10, E11]: Fit Wave/Klein-Gordon/Duffing/Sine-Gordon and exponential/stretched/bi-exponential/Mittag-Leffler families on discovery blocks and evaluate BIC.
\item \textbf{S5} [G4; equations E6]: Evaluate frozen held-out temporal blocks with NRMSE ratios.
\item \textbf{S6} [G5; equations E12, E13]: Apply deterministic coarsening diagnostics and strict C++/OpenMP FP64 parity to confirmed candidates.
\item \textbf{S7} [G6; equations E12]: Consume optional canonical-framework SYCL/DD32 smoke evidence as screening-only diagnostics.
\item \textbf{S8} [G7; equations E5, E6]: Assess evidence-class-aware replication; synthetic/simulation controls cannot close physical cross-source replication.
\item \textbf{S9} [G8; equations E4, E8]: Classify joint, phase-only or memory-only physical support only after eligible replication.
\end{enumerate}
}{}

% SST-REPORT-SECTION:SOURCES
\section{Sources and provider contract}
Real inputs satisfy \code{A056-DYNAMIC-PROVIDER-5}. Static knot geometry is not a substitute for $\varphi(t,s)$. E010 supplies qualified geometry/provenance; A056 supplies the perturbation experiment and dynamics. Each physical case must carry a frozen phase definition, perturbation ID, solver/provider version, carrier and upstream geometry hashes, evidence class, independence unit and provenance family. The bundled synthetic cases are implementation controls only.

\subsection{Real-provider construction}
The frozen Kelvin perturbation and observables are
\begin{align}
\mathbf X_\epsilon(s,0)&=\mathbf X_0(s)+\epsilon[\cos(m\theta)\mathbf e_1+\sin(m\theta)\mathbf e_2],\\
\dot{\mathbf X}(s,t)&=\mathbf u(\mathbf X(s,t),t),\\
q(s,t)&=\delta\mathbf X\cdot\mathbf e_1+i\,\delta\mathbf X\cdot\mathbf e_2,\\
\varphi(s,t)&=\arg[q(s,t)e^{-im\theta(s)}],\\
R(t)&=\frac{|c_m(t)|}{|c_m(0)|},\qquad c_m=N_s^{-1}\sum_jq(s_j,t)e^{-im\theta_j}.
\end{align}
The volumetric lane evolves the 3-D incompressible, unforced Euler equations with a 2/3-dealiased pseudo-spectral RK4 method. Its finite-core seed is a Gaussian tangent tube,
\[
\boldsymbol\omega(\mathbf x,0)\propto\sum_j\mathbf T_j\exp[-d_{\rm per}(\mathbf x,\mathbf X_j)^2/(2\sigma^2)].
\]
The cheaper filament lane evolves the regularized Biot--Savart centerline directly. Both lanes are classified as simulation evidence.

% SST-REPORT-SECTION:GATES
\section{Gates}
\begin{description}
\item[G0] frozen protocol and blind commitment;
\item[G1] provider/admissibility;
\item[G2] Python-FP64 analytic and spectral qualification;
\item[G3] blind discovery;
\item[G4] held-out confirmation;
\item[G5] deterministic resolution checks and C++/OpenMP FP64 certification;
\item[G6] optional SYCL FP32/DD32 screening, non-blocking for physical replication;
\item[G7] evidence-class-aware physical cross-source replication;
\item[G8] joint mechanism classification;
\item[G9] commitment-verified reveal.
\end{description}
Synthetic/simulation evidence cannot close G7.

% SST-REPORT-SECTION:NUMERICS
\section{Numerics and convergence}
Phase support requires positive Sine--Gordon coefficients, condition number $\le 10^8$, discovery $\Delta\mathrm{BIC}\ge10$, the registered held-out NRMSE-ratio threshold, and the registered absolute held-out NRMSE ceiling. Provider/scoring configuration SHA-256 values are embedded in the frozen science contract. Memory support requires the same BIC/holdout logic and absolute NRMSE ceiling plus $0.48<\alpha<0.97$.

For $\mathbf c=(a,b)$,
\[
\delta_c=\frac{\|\mathbf c_f-\mathbf c_c\|_2}{\max(\|\mathbf c_f\|_2,10^{-30})}.
\]
For memory,
\[
\delta_\alpha=|\alpha_f-\alpha_c|,\qquad
\delta_\tau=\frac{|\tau_f-\tau_c|}{\max(|\tau_f|,10^{-30})}.
\]
Production profiles additionally require A056 C++ FP64 parity $\le10^{-10}$.

% SST-REPORT-SECTION:BACKENDS
\section{Backend authority}
\begin{longtable}{lll}\toprule Lane & Precision & Authority\\\midrule
Python/NumPy & FP64 & REFERENCE\\
C++17/OpenMP & FP64 & CERTIFICATION\\
SYCL & FP32 & SCREENING\_ONLY\\
SYCL DD32 / FP32x2 & two FP32 components & SCREENING\_ONLY\\\bottomrule\end{longtable}
DD32 is not IEEE FP64; its accuracy is measured against the FP64 reference. The canonical v1.0.4 external worker and oneAPI link/provenance policy are used unchanged.

% SST-REPORT-SECTION:BLINDNESS
\section{Blindness and reveal}
The framework creates nonced commitments for the private forbidden-term list and reveal payload. Blind source/output packaging excludes private/revealed content. Reveal is explicit and allowed only after a frozen FULL/CERTIFY campaign satisfying the configured reveal preconditions.

% SST-REPORT-SECTION:STATISTICS
\section{Statistics and model selection}
Discovery uses
\[
\mathrm{BIC}=n\ln(\mathrm{RSS}/n)+k\ln n.
\]
Confirmation uses normalized holdout error. Sine--Gordon and Mittag--Leffler must each improve on the best registered competitor, not merely fit acceptably in isolation. No threshold is tuned after reveal.

% SST-REPORT-SECTION:RESULTS
\section{Results}
\IfFileExists{A056_SUMMARY.tex}{\input{A056_SUMMARY.tex}}{No A056 run summary has been generated in this report directory.}
\IfFileExists{A056_CASE_RESULTS.tex}{\input{A056_CASE_RESULTS.tex}}{}
\IfFileExists{AUTO_RESULTS.tex}{\input{AUTO_RESULTS.tex}}{}

% SST-REPORT-SECTION:INTERPRETATION
\section{Interpretation}
Synthetic-control recovery validates implementation/model-selection behavior only. The new E010 finite-core/Euler lane can establish a replicated simulation result, but simulation-only support remains ineligible for physical G7. Physical support requires G5 certification and G7 eligible replication. A joint SST mechanism claim requires both phase and memory closure to replicate jointly and close G8.

% SST-REPORT-SECTION:REPRODUCIBILITY
\section{Reproducibility}
The frozen protocol bundle hashes \code{falsifier.toml}, science/source/gate/blind contracts, both commitment files and this report. Run outputs record environment, source, backend, gate and output manifests. The framework is pinned in \code{FRAMEWORK\_PIN.json}; A056's instance-local C++ source is hash-bound by the canonical framework build provenance.

\begin{thebibliography}{99}
\bibitem{ScottChuMcLaughlin1973} A.~C.~Scott, F.~Y.~F.~Chu, and D.~W.~McLaughlin, ``The soliton: A new concept in applied science,'' \emph{Proceedings of the IEEE} \textbf{61} (1973), 1443--1483. \href{https://doi.org/10.1109/PROC.1973.9296}{doi:10.1109/PROC.1973.9296}.
\bibitem{Podlubny1999} I.~Podlubny, \emph{Fractional Differential Equations}, Academic Press, 1999.
\bibitem{Gorenflo2014} R.~Gorenflo, A.~A.~Kilbas, F.~Mainardi, and S.~V.~Rogosin, \emph{Mittag-Leffler Functions, Related Topics and Applications}, Springer, 2014. \href{https://doi.org/10.1007/978-3-662-43930-2}{doi:10.1007/978-3-662-43930-2}.
\bibitem{Jolliffe2002} I.~T.~Jolliffe, \emph{Principal Component Analysis}, 2nd ed., Springer, 2002. \href{https://doi.org/10.1007/b98835}{doi:10.1007/b98835}.
\bibitem{Bishop1975} R.~L.~Bishop, ``There Is More Than One Way to Frame a Curve,'' \emph{American Mathematical Monthly} \textbf{82} (1975), 246--251. \href{https://doi.org/10.2307/2319846}{doi:10.2307/2319846}.
\bibitem{Saffman1992} P.~G.~Saffman, \emph{Vortex Dynamics}, Cambridge University Press, 1992. \href{https://doi.org/10.1017/CBO9780511624063}{doi:10.1017/CBO9780511624063}.
\bibitem{Dekker1971} T.~J.~Dekker, ``A floating-point technique for extending the available precision,'' \emph{Numerische Mathematik} \textbf{18} (1971), 224--242. \href{https://doi.org/10.1007/BF01397083}{doi:10.1007/BF01397083}.
\end{thebibliography}
\end{document}
